\begin{helper}
Let $\mathcal F$ be a finite multi-hypergraph, let $f$ and $g$ be distinct
edges, and suppose that two distinct vertices $u$ and $v$ both lie in
$f\cap g$.  In the split-edge hypergraph
$\operatorname{SplitEdge}_k(\mathcal F,g,v)$ of
\noderef{KUniformSplitEdgeHypergraph}, formed by replacing $g$ with
$(g\setminus\{v\})\cup\{x_1\}$ and adding the fresh edge through $v$, the
number of independent pairs in the line-graph neighborhood of the old-edge
copy $\operatorname{inl}(f)$ is at most
\[
  P_{L(\mathcal F)}(f)+
  \bigl|N_{L(\mathcal F)}(f)\setminus\{g\}\bigr|.
\]
\end{helper}
\begin{proof}
By \noderef{IndependentPairCount}, an independent pair around an edge is a
two-element subset of its line-graph neighborhood whose two members are not
adjacent to each other.  By \noderef{LineGraphOfHypergraph}, line-graph
adjacency is precisely nonempty intersection of the corresponding hyperedges,
with distinct labels.  By \noderef{KUniformSplitEdgeHypergraph}, every old edge
other than $g$ is copied unchanged, the old label $g$ is replaced by
$(g\setminus\{v\})\cup\{x_1\}$, and one new edge $\{v,x_1,\ldots,x_{k-1}\}$ is
added.  Since $u\in f\cap g$ and $u\ne v$, the modified old copy of $g$ still
meets the old copy of $f$.

Partition the split independent pairs around $\operatorname{inl}(f)$ into
pairs not using the new edge and pairs using the new edge.  Among the first
class, send a pair of old-edge copies to the corresponding pair of old labels.
If that old pair was already independent in $L(\mathcal F)$, it is counted by
$P_{L(\mathcal F)}(f)$.  Otherwise its old adjacency must have been destroyed
by deleting $v$ from $g$: intersections among two old labels different from
$g$ are unchanged by the split, and intersections of $f$ with any old neighbor
are preserved, while $g$ can lose only the vertex $v$.  Thus every exceptional
first-class pair has the form $\{\operatorname{inl}(g),\operatorname{inl}(h)\}$
with $h\in N_{L(\mathcal F)}(f)\setminus\{g\}$ and $v\in h$.

For a second-class pair, the pair has the form
$\{\text{new},\operatorname{inl}(h)\}$.  The label $h$ is an old neighbor of
$f$, and $h\ne g$, because the new edge meets the modified old copy of $g$ at
$x_1$.  The pair is independent exactly when the new edge does not meet the
old copy of $h$, which for $h\ne g$ means $v\notin h$.  Hence these pairs
inject into the complementary subset of
$N_{L(\mathcal F)}(f)\setminus\{g\}$ consisting of neighbors not containing
$v$.

The ordinary old independent pairs, the exceptional pairs charged to neighbors
containing $v$, and the new-edge pairs charged to neighbors not containing
$v$ together exhaust the split independent-pair count.  The last two charge
sets are disjoint subsets of $N_{L(\mathcal F)}(f)\setminus\{g\}$, so adding
the three bounds gives the claimed inequality.
\end{proof}
